% =================================================
% MODELO DE RESUMO EXPANDIDO - PEPICT / UNIFESP
% =================================================

\documentclass[12pt,a4paper]{article}

% -------------------------------------------------
% PACOTES
% -------------------------------------------------

\usepackage[utf8]{inputenc}
\usepackage[T1]{fontenc}
\usepackage[brazil]{babel}
\usepackage{lmodern}
\usepackage{geometry}
\usepackage{setspace}
\usepackage{titlesec}

% -------------------------------------------------
% CONFIGURAÇÕES ABNT
% -------------------------------------------------

\geometry{
    left=3cm,
    right=2cm,
    top=3cm,
    bottom=2cm
}

\setlength{\parindent}{1.25cm}

\onehalfspacing

% -------------------------------------------------
% INÍCIO
% -------------------------------------------------

\begin{document}

% -------------------------------------------------
% CABEÇALHO INSTITUCIONAL
% -------------------------------------------------

\begin{center}

{\large \textbf{UNIVERSIDADE FEDERAL DE SÃO PAULO -- UNIFESP}}

\vspace{0.3cm}

{\normalsize PROGRAMA PEPICT}

\vspace{0.3cm}

{\normalsize ICT -- Instituto de Ciência e Tecnologia}

\vspace{1cm}

{\Large \textbf{Projeto - Panda}}

\end{center}

\vspace{1cm}

% -------------------------------------------------
% AUTORES
% -------------------------------------------------

\noindent
\textbf{Autores(as)}

Ana Carolina

\vspace{0.7cm}

% -------------------------------------------------
% DISCIPLINA E CURSO
% -------------------------------------------------

\noindent
\textbf{Disciplina e Curso}

Arquitetura e Organização de Computadores

\vspace{0.7cm}

% -------------------------------------------------
% PROGRAMA PEPICT
% -------------------------------------------------

\noindent
\textbf{Programa PEPICT vinculado}

Tecnologia e Inovação para o Desenvolvimento Social

\vspace{0.7cm}

% -------------------------------------------------
% OBJETIVOS
% -------------------------------------------------

\section*{Objetivos do Trabalho}

muitos

% -------------------------------------------------
% METODOLOGIA
% -------------------------------------------------

\section*{Metodologia}

o que eu quiser

% -------------------------------------------------
% RESULTADOS
% -------------------------------------------------

\section*{Resultados e Produtos}

funcionou

% -------------------------------------------------
% ODS
% -------------------------------------------------

\section*{Relação com os ODS}

ODS 1 - Erradicação da Pobreza

% -------------------------------------------------
% REFLEXÃO CRÍTICA
% -------------------------------------------------

\section*{Reflexão Crítica dos Estudantes}

refletido

% -------------------------------------------------
% REFERÊNCIAS
% -------------------------------------------------

\section*{Referências}

ual

\end{document}